\documentclass[11pt,a4paper]{letter}
\usepackage[margin=2.5cm]{geometry}
\usepackage{xcolor}
\newcommand{\todo}[1]{\textcolor{red}{[TODO: #1]}}
\signature{Samarth B N\\Department of Electronics and Communication Engineering\\RV College of Engineering, Bengaluru 560059, India\\samarthbn.ec25@rvce.edu.in}
\address{}
\begin{document}
\begin{letter}{Editor-in-Chief\\Advances in Space Research}
\opening{Dear Editor,}

I submit the manuscript ``Physics-guided versus data-driven forecasts of relativistic electrons at
geostationary orbit: skill, robustness and transfer across four satellites'' for consideration as a research
article in \emph{Advances in Space Research}.

Forecasts of \(>2\)\,MeV electrons at geostationary orbit are needed to protect satellites from internal
charging, including at longitudes that no GOES satellite monitors. The manuscript compares a physics-guided
forecast network with 22 other model families at lead times of 1--24\,h, using a leakage-controlled protocol, ten
seeds on three disjoint test periods and paired significance tests. It then tests transfer from GOES-15 to
GOES-16, GOES-18 and India's GSAT-19 GRASP. The results show that the physics-inspired components do not change
accuracy significantly, that a plain LSTM forecasting the change from persistence matches the full physics-guided
model and transfers best to other satellites, and that predicting changes rather than flux levels is what survives
a calibration and longitude offset. I believe this fits the journal's scope in space weather and the space environment.

\textbf{Prior version.} An earlier, shorter version of this work (\emph{STORM-PhysNet: Multi-Horizon Forecasting
of Relativistic Electron Hazards to Geostationary Satellites with Transfer to GSAT-19}) was accepted at the
AGERS 2026 conference and withdrawn before publication; it has not been published or
posted elsewhere. The present manuscript extends it with a 24\,h lead time, fourteen additional model families,
probabilistic evaluation, storm-time and alert scores, and transfer to GOES-16 and GOES-18.
\todo{confirm the withdrawal date and that no preprint exists}

The manuscript is not under consideration elsewhere. The code, data processing and predictions are openly
available. I suggest Section EM of the journal.

\closing{Yours sincerely,}
\end{letter}
\end{document}
